% !TeX root = ../../Book1.tex
\section{局部化}
\label{b1:sec:2204}

前三节反复使用“先乘截断，再磨光或延拓”。现在把其中的估计单独整理出来。局部化不仅用来拼接函数，也能把一个只允许紧支集测试函数的积分恒等式，转成较小区域上的定量控制。需要留意的是截断函数的导数：它们记录了较小区域与外部边界之间的距离，也是局部估计中常数的来源。

\subsection{截断函数}
\label{b1:subsec:220401}

设 \(U,V\subseteq\Omega\) 为开集，且 \(\overline U\Subset V\)。由光滑截断引理，可选 \(0\le\eta\le1\)，使 \(\eta\in C_c^\infty(V)\)，并在 \(\overline U\) 的邻域内等于 \(1\)。对 \(u\in W^{k,p}_{\mathrm{loc}}(\Omega)\)，函数 \(\eta\cdot u\) 在 \(U\) 上与 \(u\) 相同，却具有内部紧支集。这个替换保留了所关注区域中的函数和全部弱导数。

\begin{proposition}\label{b1:prop:sobolev-cutoff-estimate}
设 \(k\in\mathbb N_0\)、\(1\le p\le\infty\)，且 \(\eta\in C_c^\infty(V)\)。对 \(u\in W^{k,p}(V)\)，把 \(\eta\cdot u\) 在 \(V\) 外补零，得到
\[
 \|\widetilde{\eta\cdot u}\|_{W^{k,p}(\mathbb R^d)}
 \le C_{d,k}
 \left(\max_{|\beta|\le k}\|\partial^\beta\eta\|_\infty\right)
 \|u\|_{W^{k,p}(V)}.
\]
更精细地，对每个 \(|\alpha|\le k\)，
\begin{equation}\label{b1:eq:sobolev-cutoff-components}
 \|\partial^\alpha(\eta\cdot u)\|_{L^p(V)}
 \le\sum_{\beta\le\alpha}\binom\alpha\beta
     \|\partial^\beta\eta\|_\infty
     \cdot\|\partial^{\alpha-\beta}u\|_{L^p(V)}.
\end{equation}
\end{proposition}
\begin{proof}
弱 Leibniz 公式给出
\[
 \partial^\alpha(\eta\cdot u)
 =\sum_{\beta\le\alpha}\binom\alpha\beta
       (\partial^\beta\eta)\cdot(\partial^{\alpha-\beta}u).
\]
对每一项使用有界函数乘法的 \(L^p\) 估计，再用三角不等式，即得分量估计。这对 \(p=\infty\) 也成立。多重指标的数目和二项式系数仅依赖于 \(d,k\)，有限维范数比较于是给出总范数估计；所用比较常数可统一取在 \(1\le p\le\infty\) 上有界。\cref{b1:lem:compact-sobolev-zero-extension}保证补零后的导数不增加边界项，也不改变范数。
\end{proof}

在两个同心球之间，截断常数可以写得更具体。若 \(0<r<R\)，存在 \(\eta\in C_c^\infty(B(x_0,R))\)，使 \(0\le\eta\le1\)、\(\eta=1\) 于 \(\overline B(x_0,r)\)，并且
\begin{equation}\label{b1:eq:ball-cutoff-derivatives}
 \|\partial^\beta\eta\|_\infty
 \le C_{d,\beta}(R-r)^{-|\beta|}
 \quad(|\beta|\ge1).
\end{equation}
例如取 \(\delta=(R-r)/4\)，将 \(B(x_0,r+2\delta)\) 的示性函数与 \(\rho_\delta\) 卷积。卷积在 \(\overline B(x_0,r)\) 附近等于 \(1\)，支集包含于 \(\overline B(x_0,r+3\delta)\subset B(x_0,R)\)。将导数放在核上，并用
\[
 \|\partial^\beta\rho_\delta\|_{L^1}
 =\delta^{-|\beta|}\|\partial^\beta\rho\|_{L^1},
\]
就得到上述界。

\ProseParagraphBreak

因此，局部化不是无代价的限制运算。例如一阶时，
\[
 \|\nabla(\eta\cdot u)\|_{L^p}
 \le \|\eta\cdot\nabla u\|_{L^p}
      +\|u\cdot\nabla\eta\|_{L^p}.
\]
第二项含有函数本身，而不只含梯度；当 \(R-r\) 缩小时，它通常会增大。把一个全域估计用于 \(\eta\cdot u\) 后，必须保留这项，不能因 \(\eta=1\) 于内球就将它删去。

\subsection{单位分解}
\label{b1:subsec:220402}

截断处理一个区域；单位分解则使相容的局部函数成为一个整体。这里先区分“局部属于 Sobolev 空间”和“具有有限全域范数”。

\begin{proposition}\label{b1:prop:sobolev-local-gluing}
设 \(k\in\mathbb N_0\)、\(1\le p\le\infty\)，且 \((U_j)_{j\in J}\) 为开集 \(\Omega\) 的开覆盖。若 \(u\in L^1_{\mathrm{loc}}(\Omega)\)，且每个限制 \(u|_{U_j}\) 都属于 \(W^{k,p}_{\mathrm{loc}}(U_j)\)，则 \(u\in W^{k,p}_{\mathrm{loc}}(\Omega)\)。

\ProseParagraphBreak

若进一步 \(u\in L^p(\Omega)\)，且各个局部弱导数拼成的函数 \(g_\alpha\) 满足 \(g_\alpha\in L^p(\Omega)\)（\(1\le|\alpha|\le k\)），则 \(u\in W^{k,p}(\Omega)\)，并有 \(\partial^\alpha u=g_\alpha\)。
\end{proposition}
\begin{proof}
在两个开集的交上，同阶弱导数由唯一性相等。因此各局部导数可以拼接。为避免不可数个零测集的问题，先利用欧氏空间的可数基，从原覆盖中取出可数子覆盖；再在这个可数族内去掉各交集上导数不一致的零测集。由此得到几乎处处定义且可测的 \(g_\alpha\)。

任取 \(K\Subset\Omega\)。可以用有限个相对紧于某些 \(U_j\) 的小开集覆盖 \(K\)，所以 \(g_\alpha\) 在 \(K\) 附近属于 \(L^p\)。现取 \(\varphi\in C_c^\infty(\Omega)\)，并将其按从属于覆盖的光滑单位分解写为有限和
\[
 \varphi=\sum_\ell\psi_\ell\varphi.
\]
每个非零分量都紧支于某个 \(U_j\)。在那里应用弱导数恒等式，求和得到
\[
 \begin{aligned}
 \int_\Omega u\,\partial^\alpha\varphi
 &=\sum_\ell\int_\Omega u\,\partial^\alpha(\psi_\ell\varphi)\\
 &=(-1)^{|\alpha|}\sum_\ell\int_\Omega g_\alpha\psi_\ell\varphi
 =(-1)^{|\alpha|}\int_\Omega g_\alpha\varphi.
 \end{aligned}
\]
第一处等号使用的是测试函数分解后的导数之和；单位分解导数产生的项已经在这个和中相消。于是 \(g_\alpha\) 是全域弱导数。局部或全域的可积性分别给出两种结论。
\end{proof}

同样，若存在 \(\Omega\) 中的局部有限闭集族 \((K_j)\)，使每个 \(v_j\in W^{k,p}_{\mathrm{loc}}(\Omega)\) 都在 \(\Omega\setminus K_j\) 上几乎处处为零，那么局部有限和
\[
 v=\sum_jv_j
\]
属于 \(W^{k,p}_{\mathrm{loc}}(\Omega)\)，其导数可以逐项相加，因为每个点附近实际只有有限项。这里没有保证 \(\|v\|_{W^{k,p}(\Omega)}<\infty\)。例如在全空间取光滑单位分解，所有分解函数都在各自的紧集外为零，但它们的和为 \(1\)，在有限 \(p\) 时不属于 \(L^p(\mathbb R^d)\)。

\ProseParagraphBreak

获得全域界需要额外估计。一种简单而充分的条件是
\[
 \sum_j\|v_j\|_{W^{k,p}(\Omega)}<\infty.
\]
此时完备性使级数在全域范数中收敛，其极限与局部有限和相同，而且总范数不超过右侧之和。Meyers--Serrin 定理正是对磨光误差建立了这样的可和界；不必要求单位分解后的原函数各项也满足它。

\subsection{局部 Sobolev 空间}
\label{b1:subsec:220403}

\secref{b1:sec:2101}的局部空间定义，要求在每个内部相对紧开集上控制函数及导数。用刚才的截断估计，可以把这个条件改写成全空间中的一族条件。

\begin{proposition}\label{b1:prop:local-sobolev-cutoff-characterization}
设 \(k\in\mathbb N_0\)、\(1\le p\le\infty\)，且 \(u\in L^1_{\mathrm{loc}}(\Omega)\)。下列条件等价：
\begin{equivalences}
\item\label{b1:item:local-sobolev-membership}
\(u\in W^{k,p}_{\mathrm{loc}}(\Omega)\)。
\item\label{b1:item:all-sobolev-cutoffs}
对每个 \(\chi\in C_c^\infty(\Omega)\)，补零函数 \(\widetilde{\chi\cdot u}\) 属于 \(W^{k,p}(\mathbb R^d)\)。
\end{equivalences}
若 \(u_n,u\in W^{k,p}_{\mathrm{loc}}(\Omega)\)，则在每个 \(\overline U\Subset\Omega\) 的开集 \(U\) 上有 \(u_n\to u\) 于 \(W^{k,p}(U)\)，当且仅当对每个上述 \(\chi\)，
\[
 \widetilde{\chi\cdot u_n}\longrightarrow\widetilde{\chi\cdot u}
 \quad\text{于 }W^{k,p}(\mathbb R^d).
\]
\end{proposition}
\begin{proof}
由局部可积性，先取一个内部开集包含 \(\operatorname{supp}\chi\)，再应用截断估计和补零引理，得到从\itemref{b1:item:local-sobolev-membership}到\itemref{b1:item:all-sobolev-cutoffs}的蕴含。反过来，对任意 \(\overline U\Subset\Omega\)，选取在 \(\overline U\) 附近为一的 \(\chi\)。由于 \(\widetilde{\chi\cdot u}|_U=u|_U\)，限制全空间弱导数便得到 \(u|_U\in W^{k,p}(U)\)。

收敛部分使用相同的两个操作。由局部收敛及固定 \(\chi\) 的截断估计，得到全空间乘积的收敛；反过来，取 \(\chi=1\) 于 \(\overline U\) 附近，限制不会增大范数，故得到 \(U\) 上的收敛。
\end{proof}

这里的“每个截断函数”不意味着实际工作中必须检验不可数多项。取递增的内部相对紧开集 \(U_j\) 穷竭 \(\Omega\)，再固定 \(\chi_j=1\) 于 \(\overline{U_j}\) 附近，即可只用这一列截断检验局部条件。任意内部紧集都包含在某个 \(U_j\) 中，因而不会漏掉局部信息。

\ProseParagraphBreak

局部有界性还足以在自反指数范围内抽取局部弱极限。具体说，若 \(1<p<\infty\)，且 \((u_n)\) 在每个 \(W^{k,p}(U_j)\) 中有界，则在 \(U_1\) 上用自反性抽取弱收敛子列，再在 \(U_2\) 上从中继续抽取，依次进行。对角子列在每个 \(U_j\) 上弱收敛。限制算子是有界线性的，故相邻区域上的极限限制相同；这些极限拼成 \(u\in W^{k,p}_{\mathrm{loc}}(\Omega)\)。这只给出局部弱收敛，不保证全域范数有界，更不保证强收敛。以互相远离的平移函数为例，函数可以在每个固定内部区域上消失，而全域范数一直保持不变。

\subsection{局部能量估计与 Caccioppoli 不等式}
\label{b1:subsec:220404}

局部化的一个重要用途，是允许测试函数依赖于待研究的 Sobolev 函数。若积分恒等式最初只对 \(C_c^\infty\) 函数成立，不能直接把一个不光滑的乘积代入；内部紧支集的范数逼近提供了所需过渡。

\ProseParagraphBreak

以下在实数范围内讨论。设 \(u\in H^1_{\mathrm{loc}}(\Omega)\)，满足
\begin{equation}\label{b1:eq:weak-harmonic-condition}
 \int_\Omega\nabla u\cdot\nabla\varphi\,\mathrm dx=0
 \quad\bigl(\varphi\in C_c^\infty(\Omega)\bigr).
\end{equation}
这个恒等式是方程 \(-\Delta u=0\) 的弱形式，并不预先要求 \(u\) 有经典二阶导数。下面的\term[caccioppoli-budengshi]{Caccioppoli 不等式}[Caccioppoli inequality]用较大球中的函数大小控制较小球中的梯度能量。

\begin{theorem}[Caccioppoli]\label{b1:thm:caccioppoli}
若 \(u\in H^1_{\mathrm{loc}}(\Omega)\) 满足\eqref{b1:eq:weak-harmonic-condition}，则对实常数 \(c\) 及实值 \(\eta\in C_c^\infty(\Omega)\)，
\begin{equation}\label{b1:eq:caccioppoli-cutoff}
 \int_\Omega\eta^2|\nabla u|^2\,\mathrm dx
 \le4\int_\Omega|u-c|^2\cdot|\nabla\eta|^2\,\mathrm dx.
\end{equation}
特别地，若 \(0<r<R\)、\(\overline B(x_0,R)\Subset\Omega\)，则
\begin{equation}\label{b1:eq:caccioppoli-balls}
 \int_{B(x_0,r)}|\nabla u|^2\,\mathrm dx
 \le\frac{C_d}{(R-r)^2}
       \int_{B(x_0,R)}|u-c|^2\,\mathrm dx.
\end{equation}
\end{theorem}
\begin{proof}
令 \(v=\eta^2(u-c)\)。弱乘积公式保证 \(v\in H^1(\Omega)\)，且它具有内部紧支集。选一个开集 \(V\)，使 \(\operatorname{supp}\eta\Subset V\)、\(\overline V\Subset\Omega\)。由\cref{b1:prop:compact-sobolev-smooth-approximation}，存在 \(v_j\in C_c^\infty(V)\)，使 \(v_j\to v\) 于 \(H^1(\Omega)\)。由于 \(\nabla u\in L^2(V)\)，Cauchy--Schwarz 不等式给出
\[
 \left|\int_V\nabla u\cdot\nabla(v_j-v)\,\mathrm dx\right|
 \le\|\nabla u\|_{L^2(V)}
      \cdot\|\nabla(v_j-v)\|_{L^2(V)}
 \longrightarrow0.
\]
因此弱恒等式也允许使用测试函数 \(v\)。展开
\[
 \nabla v=\eta^2\nabla u+2\eta(u-c)\nabla\eta
\]
并代入，得到
\[
 \int_\Omega\eta^2|\nabla u|^2
 =-2\int_\Omega\eta(u-c)\nabla u\cdot\nabla\eta.
\]
用向量 Cauchy--Schwarz 不等式及 \(2ab\le a^2/2+2b^2\)，得到
\[
 \int_\Omega\eta^2|\nabla u|^2
 \le\frac12\int_\Omega\eta^2|\nabla u|^2
       +2\int_\Omega|u-c|^2\cdot|\nabla\eta|^2.
\]
吸收左侧的一半即得\eqref{b1:eq:caccioppoli-cutoff}。最后选取\eqref{b1:eq:ball-cutoff-derivatives}中的球间截断，使用 \(\eta=1\) 于内球及 \(|\nabla\eta|\le C_d/(R-r)\)，便得\eqref{b1:eq:caccioppoli-balls}。
\end{proof}

常数 \(c\) 可以自由选取，因为方程只涉及梯度。取 \(c=0\) 得到直接的 \(L^2\) 控制；取 \(c\) 为较大球上的平均值，则右侧只测量函数相对于常数的振幅。当 \(u\) 本来就是常数时，左右两侧都为零，这也说明估计保留了方程对加常数的不变性。

\ProseParagraphBreak

本证明没有从 \(\|u\|_{L^2}\) 推出任意函数的梯度估计，而是先用方程将梯度能量变成截断交叉项。没有弱恒等式时，这一步并不存在。另一方面，截断只在内部起作用，所以定理没有使用边界值或迹；靠近边界的估计则需要额外信息。下一章将给出 Sobolev 函数的边界值应如何定义，以及这种边界值与 \(W_0^{1,p}\) 的关系。

单位分解的归一化与支撑安排见\cref{b1:fig:visual-f16}。
% F16：本书原创矢量示意；依据所在节的定义、证明或明确模型。
\begin{figure}[htbp]
\centering
\begin{tikzpicture}[x=1cm,y=1cm,>=stealth,every node/.style={font=\small},line width=.65pt]
\node at(6,3.7){开覆盖与截断函数};
\draw[AnalysisBlue,thick](1.0,2.7)--(5.0,2.7);\node[AnalysisBlue] at(3.0,3.1){\(U_1\)};
\draw[AnalysisBlue,thick] plot coordinates {(1.0000,0.0000) (1.0400,0.0000) (1.0800,0.0000) (1.1200,0.0174) (1.1600,0.0521) (1.2000,0.0868) (1.2400,0.1216) (1.2800,0.1563) (1.3200,0.1911) (1.3600,0.2258) (1.4000,0.2605) (1.4400,0.2953) (1.4800,0.3300) (1.5200,0.3647) (1.5600,0.3995) (1.6000,0.4342) (1.6400,0.4689) (1.6800,0.5037) (1.7200,0.5384) (1.7600,0.5732) (1.8000,0.6079) (1.8400,0.6426) (1.8800,0.6774) (1.9200,0.7121) (1.9600,0.7468) (2.0000,0.7816) (2.0400,0.8163) (2.0800,0.8511) (2.1200,0.8858) (2.1600,0.9205) (2.2000,0.9553) (2.2400,0.9900) (2.2800,1.0247) (2.3200,1.0595) (2.3600,1.0942) (2.4000,1.1289) (2.4400,1.1637) (2.4800,1.1984) (2.5200,1.2332) (2.5600,1.2679) (2.6000,1.3026) (2.6400,1.3374) (2.6800,1.3721) (2.7200,1.4068) (2.7600,1.4416) (2.8000,1.4763) (2.8400,1.5111) (2.8800,1.5458) (2.9200,1.5805) (2.9600,1.6153) (3.0000,1.6500) (3.0400,1.6153) (3.0800,1.5805) (3.1200,1.5458) (3.1600,1.5111) (3.2000,1.4763) (3.2400,1.4416) (3.2800,1.4068) (3.3200,1.3721) (3.3600,1.3374) (3.4000,1.3026) (3.4400,1.2679) (3.4800,1.2332) (3.5200,1.1984) (3.5600,1.1637) (3.6000,1.1289) (3.6400,1.0942) (3.6800,1.0595) (3.7200,1.0247) (3.7600,0.9900) (3.8000,0.9553) (3.8400,0.9205) (3.8800,0.8858) (3.9200,0.8511) (3.9600,0.8163) (4.0000,0.7816) (4.0400,0.7468) (4.0800,0.7121) (4.1200,0.6774) (4.1600,0.6426) (4.2000,0.6079) (4.2400,0.5732) (4.2800,0.5384) (4.3200,0.5037) (4.3600,0.4689) (4.4000,0.4342) (4.4400,0.3995) (4.4800,0.3647) (4.5200,0.3300) (4.5600,0.2953) (4.6000,0.2605) (4.6400,0.2258) (4.6800,0.1911) (4.7200,0.1563) (4.7600,0.1216) (4.8000,0.0868) (4.8400,0.0521) (4.8800,0.0174) (4.9200,0.0000) (4.9600,0.0000) (5.0000,0.0000)};
\draw[orange!80!black,thick](4.5,2.7)--(8.5,2.7);\node[orange!80!black] at(6.5,3.1){\(U_2\)};
\draw[orange!80!black,thick] plot coordinates {(4.5000,0.0000) (4.5400,0.0000) (4.5800,0.0000) (4.6200,0.0174) (4.6600,0.0521) (4.7000,0.0868) (4.7400,0.1216) (4.7800,0.1563) (4.8200,0.1911) (4.8600,0.2258) (4.9000,0.2605) (4.9400,0.2953) (4.9800,0.3300) (5.0200,0.3647) (5.0600,0.3995) (5.1000,0.4342) (5.1400,0.4689) (5.1800,0.5037) (5.2200,0.5384) (5.2600,0.5732) (5.3000,0.6079) (5.3400,0.6426) (5.3800,0.6774) (5.4200,0.7121) (5.4600,0.7468) (5.5000,0.7816) (5.5400,0.8163) (5.5800,0.8511) (5.6200,0.8858) (5.6600,0.9205) (5.7000,0.9553) (5.7400,0.9900) (5.7800,1.0247) (5.8200,1.0595) (5.8600,1.0942) (5.9000,1.1289) (5.9400,1.1637) (5.9800,1.1984) (6.0200,1.2332) (6.0600,1.2679) (6.1000,1.3026) (6.1400,1.3374) (6.1800,1.3721) (6.2200,1.4068) (6.2600,1.4416) (6.3000,1.4763) (6.3400,1.5111) (6.3800,1.5458) (6.4200,1.5805) (6.4600,1.6153) (6.5000,1.6500) (6.5400,1.6153) (6.5800,1.5805) (6.6200,1.5458) (6.6600,1.5111) (6.7000,1.4763) (6.7400,1.4416) (6.7800,1.4068) (6.8200,1.3721) (6.8600,1.3374) (6.9000,1.3026) (6.9400,1.2679) (6.9800,1.2332) (7.0200,1.1984) (7.0600,1.1637) (7.1000,1.1289) (7.1400,1.0942) (7.1800,1.0595) (7.2200,1.0247) (7.2600,0.9900) (7.3000,0.9553) (7.3400,0.9205) (7.3800,0.8858) (7.4200,0.8511) (7.4600,0.8163) (7.5000,0.7816) (7.5400,0.7468) (7.5800,0.7121) (7.6200,0.6774) (7.6600,0.6426) (7.7000,0.6079) (7.7400,0.5732) (7.7800,0.5384) (7.8200,0.5037) (7.8600,0.4689) (7.9000,0.4342) (7.9400,0.3995) (7.9800,0.3647) (8.0200,0.3300) (8.0600,0.2953) (8.1000,0.2605) (8.1400,0.2258) (8.1800,0.1911) (8.2200,0.1563) (8.2600,0.1216) (8.3000,0.0868) (8.3400,0.0521) (8.3800,0.0174) (8.4200,0.0000) (8.4600,0.0000) (8.5000,0.0000)};
\draw[black,thick](8.0,2.7)--(12.0,2.7);\node[black] at(10.0,3.1){\(U_3\)};
\draw[black,thick] plot coordinates {(8.0000,0.0000) (8.0400,0.0000) (8.0800,0.0000) (8.1200,0.0174) (8.1600,0.0521) (8.2000,0.0868) (8.2400,0.1216) (8.2800,0.1563) (8.3200,0.1911) (8.3600,0.2258) (8.4000,0.2605) (8.4400,0.2953) (8.4800,0.3300) (8.5200,0.3647) (8.5600,0.3995) (8.6000,0.4342) (8.6400,0.4689) (8.6800,0.5037) (8.7200,0.5384) (8.7600,0.5732) (8.8000,0.6079) (8.8400,0.6426) (8.8800,0.6774) (8.9200,0.7121) (8.9600,0.7468) (9.0000,0.7816) (9.0400,0.8163) (9.0800,0.8511) (9.1200,0.8858) (9.1600,0.9205) (9.2000,0.9553) (9.2400,0.9900) (9.2800,1.0247) (9.3200,1.0595) (9.3600,1.0942) (9.4000,1.1289) (9.4400,1.1637) (9.4800,1.1984) (9.5200,1.2332) (9.5600,1.2679) (9.6000,1.3026) (9.6400,1.3374) (9.6800,1.3721) (9.7200,1.4068) (9.7600,1.4416) (9.8000,1.4763) (9.8400,1.5111) (9.8800,1.5458) (9.9200,1.5805) (9.9600,1.6153) (10.0000,1.6500) (10.0400,1.6153) (10.0800,1.5805) (10.1200,1.5458) (10.1600,1.5111) (10.2000,1.4763) (10.2400,1.4416) (10.2800,1.4068) (10.3200,1.3721) (10.3600,1.3374) (10.4000,1.3026) (10.4400,1.2679) (10.4800,1.2332) (10.5200,1.1984) (10.5600,1.1637) (10.6000,1.1289) (10.6400,1.0942) (10.6800,1.0595) (10.7200,1.0247) (10.7600,0.9900) (10.8000,0.9553) (10.8400,0.9205) (10.8800,0.8858) (10.9200,0.8511) (10.9600,0.8163) (11.0000,0.7816) (11.0400,0.7468) (11.0800,0.7121) (11.1200,0.6774) (11.1600,0.6426) (11.2000,0.6079) (11.2400,0.5732) (11.2800,0.5384) (11.3200,0.5037) (11.3600,0.4689) (11.4000,0.4342) (11.4400,0.3995) (11.4800,0.3647) (11.5200,0.3300) (11.5600,0.2953) (11.6000,0.2605) (11.6400,0.2258) (11.6800,0.1911) (11.7200,0.1563) (11.7600,0.1216) (11.8000,0.0868) (11.8400,0.0521) (11.8800,0.0174) (11.9200,0.0000) (11.9600,0.0000) (12.0000,0.0000)};
\node at(6,-.65){\(\psi_j=\eta_j/\sum_i\eta_i\)，在工作区域上 \(\sum_j\psi_j=1\)};
\draw[AnalysisBlue,thick] plot coordinates {(1.2000,-1.3000) (1.2463,-1.3000) (1.2926,-1.3000) (1.3389,-1.3000) (1.3852,-1.3000) (1.4314,-1.3000) (1.4777,-1.3000) (1.5240,-1.3000) (1.5703,-1.3000) (1.6166,-1.3000) (1.6629,-1.3000) (1.7092,-1.3000) (1.7555,-1.3000) (1.8017,-1.3000) (1.8480,-1.3000) (1.8943,-1.3000) (1.9406,-1.3000) (1.9869,-1.3000) (2.0332,-1.3000) (2.0795,-1.3000) (2.1258,-1.3000) (2.1721,-1.3000) (2.2183,-1.3000) (2.2646,-1.3000) (2.3109,-1.3000) (2.3572,-1.3000) (2.4035,-1.3000) (2.4498,-1.3000) (2.4961,-1.3000) (2.5424,-1.3000) (2.5886,-1.3000) (2.6349,-1.3000) (2.6812,-1.3000) (2.7275,-1.3000) (2.7738,-1.3000) (2.8201,-1.3000) (2.8664,-1.3000) (2.9127,-1.3000) (2.9590,-1.3000) (3.0052,-1.3000) (3.0515,-1.3000) (3.0978,-1.3000) (3.1441,-1.3000) (3.1904,-1.3000) (3.2367,-1.3000) (3.2830,-1.3000) (3.3293,-1.3000) (3.3755,-1.3000) (3.4218,-1.3000) (3.4681,-1.3000) (3.5144,-1.3000) (3.5607,-1.3000) (3.6070,-1.3000) (3.6533,-1.3000) (3.6996,-1.3000) (3.7459,-1.3000) (3.7921,-1.3000) (3.8384,-1.3000) (3.8847,-1.3000) (3.9310,-1.3000) (3.9773,-1.3000) (4.0236,-1.3000) (4.0699,-1.3000) (4.1162,-1.3000) (4.1624,-1.3000) (4.2087,-1.3000) (4.2550,-1.3000) (4.3013,-1.3000) (4.3476,-1.3000) (4.3939,-1.3000) (4.4402,-1.3000) (4.4865,-1.3000) (4.5328,-1.3000) (4.5790,-1.3000) (4.6253,-1.4266) (4.6716,-1.6581) (4.7179,-1.8895) (4.7642,-2.1210) (4.8105,-2.3524) (4.8568,-2.5838) (4.9031,-2.8000) (4.9493,-2.8000) (4.9956,-2.8000) (5.0419,-2.8000) (5.0882,-2.8000) (5.1345,-2.8000) (5.1808,-2.8000) (5.2271,-2.8000) (5.2734,-2.8000) (5.3197,-2.8000) (5.3659,-2.8000) (5.4122,-2.8000) (5.4585,-2.8000) (5.5048,-2.8000) (5.5511,-2.8000) (5.5974,-2.8000) (5.6437,-2.8000) (5.6900,-2.8000) (5.7362,-2.8000) (5.7825,-2.8000) (5.8288,-2.8000) (5.8751,-2.8000) (5.9214,-2.8000) (5.9677,-2.8000) (6.0140,-2.8000) (6.0603,-2.8000) (6.1066,-2.8000) (6.1528,-2.8000) (6.1991,-2.8000) (6.2454,-2.8000) (6.2917,-2.8000) (6.3380,-2.8000) (6.3843,-2.8000) (6.4306,-2.8000) (6.4769,-2.8000) (6.5231,-2.8000) (6.5694,-2.8000) (6.6157,-2.8000) (6.6620,-2.8000) (6.7083,-2.8000) (6.7546,-2.8000) (6.8009,-2.8000) (6.8472,-2.8000) (6.8934,-2.8000) (6.9397,-2.8000) (6.9860,-2.8000) (7.0323,-2.8000) (7.0786,-2.8000) (7.1249,-2.8000) (7.1712,-2.8000) (7.2175,-2.8000) (7.2638,-2.8000) (7.3100,-2.8000) (7.3563,-2.8000) (7.4026,-2.8000) (7.4489,-2.8000) (7.4952,-2.8000) (7.5415,-2.8000) (7.5878,-2.8000) (7.6341,-2.8000) (7.6803,-2.8000) (7.7266,-2.8000) (7.7729,-2.8000) (7.8192,-2.8000) (7.8655,-2.8000) (7.9118,-2.8000) (7.9581,-2.8000) (8.0044,-2.8000) (8.0507,-2.8000) (8.0969,-2.8000) (8.1432,-2.8000) (8.1895,-2.8000) (8.2358,-2.8000) (8.2821,-2.8000) (8.3284,-2.8000) (8.3747,-2.8000) (8.4210,-2.8000) (8.4672,-2.8000) (8.5135,-2.8000) (8.5598,-2.8000) (8.6061,-2.8000) (8.6524,-2.8000) (8.6987,-2.8000) (8.7450,-2.8000) (8.7913,-2.8000) (8.8376,-2.8000) (8.8838,-2.8000) (8.9301,-2.8000) (8.9764,-2.8000) (9.0227,-2.8000) (9.0690,-2.8000) (9.1153,-2.8000) (9.1616,-2.8000) (9.2079,-2.8000) (9.2541,-2.8000) (9.3004,-2.8000) (9.3467,-2.8000) (9.3930,-2.8000) (9.4393,-2.8000) (9.4856,-2.8000) (9.5319,-2.8000) (9.5782,-2.8000) (9.6245,-2.8000) (9.6707,-2.8000) (9.7170,-2.8000) (9.7633,-2.8000) (9.8096,-2.8000) (9.8559,-2.8000) (9.9022,-2.8000) (9.9485,-2.8000) (9.9948,-2.8000) (10.0410,-2.8000) (10.0873,-2.8000) (10.1336,-2.8000) (10.1799,-2.8000) (10.2262,-2.8000) (10.2725,-2.8000) (10.3188,-2.8000) (10.3651,-2.8000) (10.4114,-2.8000) (10.4576,-2.8000) (10.5039,-2.8000) (10.5502,-2.8000) (10.5965,-2.8000) (10.6428,-2.8000) (10.6891,-2.8000) (10.7354,-2.8000) (10.7817,-2.8000) (10.8279,-2.8000) (10.8742,-2.8000) (10.9205,-2.8000) (10.9668,-2.8000) (11.0131,-2.8000) (11.0594,-2.8000) (11.1057,-2.8000) (11.1520,-2.8000) (11.1983,-2.8000) (11.2445,-2.8000) (11.2908,-2.8000) (11.3371,-2.8000) (11.3834,-2.8000) (11.4297,-2.8000) (11.4760,-2.8000) (11.5223,-2.8000) (11.5686,-2.8000) (11.6148,-2.8000) (11.6611,-2.8000) (11.7074,-2.8000) (11.7537,-2.8000) (11.8000,-2.8000)};
\draw[orange!80!black,thick] plot coordinates {(1.2000,-2.8000) (1.2463,-2.8000) (1.2926,-2.8000) (1.3389,-2.8000) (1.3852,-2.8000) (1.4314,-2.8000) (1.4777,-2.8000) (1.5240,-2.8000) (1.5703,-2.8000) (1.6166,-2.8000) (1.6629,-2.8000) (1.7092,-2.8000) (1.7555,-2.8000) (1.8017,-2.8000) (1.8480,-2.8000) (1.8943,-2.8000) (1.9406,-2.8000) (1.9869,-2.8000) (2.0332,-2.8000) (2.0795,-2.8000) (2.1258,-2.8000) (2.1721,-2.8000) (2.2183,-2.8000) (2.2646,-2.8000) (2.3109,-2.8000) (2.3572,-2.8000) (2.4035,-2.8000) (2.4498,-2.8000) (2.4961,-2.8000) (2.5424,-2.8000) (2.5886,-2.8000) (2.6349,-2.8000) (2.6812,-2.8000) (2.7275,-2.8000) (2.7738,-2.8000) (2.8201,-2.8000) (2.8664,-2.8000) (2.9127,-2.8000) (2.9590,-2.8000) (3.0052,-2.8000) (3.0515,-2.8000) (3.0978,-2.8000) (3.1441,-2.8000) (3.1904,-2.8000) (3.2367,-2.8000) (3.2830,-2.8000) (3.3293,-2.8000) (3.3755,-2.8000) (3.4218,-2.8000) (3.4681,-2.8000) (3.5144,-2.8000) (3.5607,-2.8000) (3.6070,-2.8000) (3.6533,-2.8000) (3.6996,-2.8000) (3.7459,-2.8000) (3.7921,-2.8000) (3.8384,-2.8000) (3.8847,-2.8000) (3.9310,-2.8000) (3.9773,-2.8000) (4.0236,-2.8000) (4.0699,-2.8000) (4.1162,-2.8000) (4.1624,-2.8000) (4.2087,-2.8000) (4.2550,-2.8000) (4.3013,-2.8000) (4.3476,-2.8000) (4.3939,-2.8000) (4.4402,-2.8000) (4.4865,-2.8000) (4.5328,-2.8000) (4.5790,-2.8000) (4.6253,-2.6734) (4.6716,-2.4419) (4.7179,-2.2105) (4.7642,-1.9790) (4.8105,-1.7476) (4.8568,-1.5162) (4.9031,-1.3000) (4.9493,-1.3000) (4.9956,-1.3000) (5.0419,-1.3000) (5.0882,-1.3000) (5.1345,-1.3000) (5.1808,-1.3000) (5.2271,-1.3000) (5.2734,-1.3000) (5.3197,-1.3000) (5.3659,-1.3000) (5.4122,-1.3000) (5.4585,-1.3000) (5.5048,-1.3000) (5.5511,-1.3000) (5.5974,-1.3000) (5.6437,-1.3000) (5.6900,-1.3000) (5.7362,-1.3000) (5.7825,-1.3000) (5.8288,-1.3000) (5.8751,-1.3000) (5.9214,-1.3000) (5.9677,-1.3000) (6.0140,-1.3000) (6.0603,-1.3000) (6.1066,-1.3000) (6.1528,-1.3000) (6.1991,-1.3000) (6.2454,-1.3000) (6.2917,-1.3000) (6.3380,-1.3000) (6.3843,-1.3000) (6.4306,-1.3000) (6.4769,-1.3000) (6.5231,-1.3000) (6.5694,-1.3000) (6.6157,-1.3000) (6.6620,-1.3000) (6.7083,-1.3000) (6.7546,-1.3000) (6.8009,-1.3000) (6.8472,-1.3000) (6.8934,-1.3000) (6.9397,-1.3000) (6.9860,-1.3000) (7.0323,-1.3000) (7.0786,-1.3000) (7.1249,-1.3000) (7.1712,-1.3000) (7.2175,-1.3000) (7.2638,-1.3000) (7.3100,-1.3000) (7.3563,-1.3000) (7.4026,-1.3000) (7.4489,-1.3000) (7.4952,-1.3000) (7.5415,-1.3000) (7.5878,-1.3000) (7.6341,-1.3000) (7.6803,-1.3000) (7.7266,-1.3000) (7.7729,-1.3000) (7.8192,-1.3000) (7.8655,-1.3000) (7.9118,-1.3000) (7.9581,-1.3000) (8.0044,-1.3000) (8.0507,-1.3000) (8.0969,-1.3000) (8.1432,-1.5162) (8.1895,-1.7476) (8.2358,-1.9790) (8.2821,-2.2105) (8.3284,-2.4419) (8.3747,-2.6734) (8.4210,-2.8000) (8.4672,-2.8000) (8.5135,-2.8000) (8.5598,-2.8000) (8.6061,-2.8000) (8.6524,-2.8000) (8.6987,-2.8000) (8.7450,-2.8000) (8.7913,-2.8000) (8.8376,-2.8000) (8.8838,-2.8000) (8.9301,-2.8000) (8.9764,-2.8000) (9.0227,-2.8000) (9.0690,-2.8000) (9.1153,-2.8000) (9.1616,-2.8000) (9.2079,-2.8000) (9.2541,-2.8000) (9.3004,-2.8000) (9.3467,-2.8000) (9.3930,-2.8000) (9.4393,-2.8000) (9.4856,-2.8000) (9.5319,-2.8000) (9.5782,-2.8000) (9.6245,-2.8000) (9.6707,-2.8000) (9.7170,-2.8000) (9.7633,-2.8000) (9.8096,-2.8000) (9.8559,-2.8000) (9.9022,-2.8000) (9.9485,-2.8000) (9.9948,-2.8000) (10.0410,-2.8000) (10.0873,-2.8000) (10.1336,-2.8000) (10.1799,-2.8000) (10.2262,-2.8000) (10.2725,-2.8000) (10.3188,-2.8000) (10.3651,-2.8000) (10.4114,-2.8000) (10.4576,-2.8000) (10.5039,-2.8000) (10.5502,-2.8000) (10.5965,-2.8000) (10.6428,-2.8000) (10.6891,-2.8000) (10.7354,-2.8000) (10.7817,-2.8000) (10.8279,-2.8000) (10.8742,-2.8000) (10.9205,-2.8000) (10.9668,-2.8000) (11.0131,-2.8000) (11.0594,-2.8000) (11.1057,-2.8000) (11.1520,-2.8000) (11.1983,-2.8000) (11.2445,-2.8000) (11.2908,-2.8000) (11.3371,-2.8000) (11.3834,-2.8000) (11.4297,-2.8000) (11.4760,-2.8000) (11.5223,-2.8000) (11.5686,-2.8000) (11.6148,-2.8000) (11.6611,-2.8000) (11.7074,-2.8000) (11.7537,-2.8000) (11.8000,-2.8000)};
\draw[black,thick] plot coordinates {(1.2000,-2.8000) (1.2463,-2.8000) (1.2926,-2.8000) (1.3389,-2.8000) (1.3852,-2.8000) (1.4314,-2.8000) (1.4777,-2.8000) (1.5240,-2.8000) (1.5703,-2.8000) (1.6166,-2.8000) (1.6629,-2.8000) (1.7092,-2.8000) (1.7555,-2.8000) (1.8017,-2.8000) (1.8480,-2.8000) (1.8943,-2.8000) (1.9406,-2.8000) (1.9869,-2.8000) (2.0332,-2.8000) (2.0795,-2.8000) (2.1258,-2.8000) (2.1721,-2.8000) (2.2183,-2.8000) (2.2646,-2.8000) (2.3109,-2.8000) (2.3572,-2.8000) (2.4035,-2.8000) (2.4498,-2.8000) (2.4961,-2.8000) (2.5424,-2.8000) (2.5886,-2.8000) (2.6349,-2.8000) (2.6812,-2.8000) (2.7275,-2.8000) (2.7738,-2.8000) (2.8201,-2.8000) (2.8664,-2.8000) (2.9127,-2.8000) (2.9590,-2.8000) (3.0052,-2.8000) (3.0515,-2.8000) (3.0978,-2.8000) (3.1441,-2.8000) (3.1904,-2.8000) (3.2367,-2.8000) (3.2830,-2.8000) (3.3293,-2.8000) (3.3755,-2.8000) (3.4218,-2.8000) (3.4681,-2.8000) (3.5144,-2.8000) (3.5607,-2.8000) (3.6070,-2.8000) (3.6533,-2.8000) (3.6996,-2.8000) (3.7459,-2.8000) (3.7921,-2.8000) (3.8384,-2.8000) (3.8847,-2.8000) (3.9310,-2.8000) (3.9773,-2.8000) (4.0236,-2.8000) (4.0699,-2.8000) (4.1162,-2.8000) (4.1624,-2.8000) (4.2087,-2.8000) (4.2550,-2.8000) (4.3013,-2.8000) (4.3476,-2.8000) (4.3939,-2.8000) (4.4402,-2.8000) (4.4865,-2.8000) (4.5328,-2.8000) (4.5790,-2.8000) (4.6253,-2.8000) (4.6716,-2.8000) (4.7179,-2.8000) (4.7642,-2.8000) (4.8105,-2.8000) (4.8568,-2.8000) (4.9031,-2.8000) (4.9493,-2.8000) (4.9956,-2.8000) (5.0419,-2.8000) (5.0882,-2.8000) (5.1345,-2.8000) (5.1808,-2.8000) (5.2271,-2.8000) (5.2734,-2.8000) (5.3197,-2.8000) (5.3659,-2.8000) (5.4122,-2.8000) (5.4585,-2.8000) (5.5048,-2.8000) (5.5511,-2.8000) (5.5974,-2.8000) (5.6437,-2.8000) (5.6900,-2.8000) (5.7362,-2.8000) (5.7825,-2.8000) (5.8288,-2.8000) (5.8751,-2.8000) (5.9214,-2.8000) (5.9677,-2.8000) (6.0140,-2.8000) (6.0603,-2.8000) (6.1066,-2.8000) (6.1528,-2.8000) (6.1991,-2.8000) (6.2454,-2.8000) (6.2917,-2.8000) (6.3380,-2.8000) (6.3843,-2.8000) (6.4306,-2.8000) (6.4769,-2.8000) (6.5231,-2.8000) (6.5694,-2.8000) (6.6157,-2.8000) (6.6620,-2.8000) (6.7083,-2.8000) (6.7546,-2.8000) (6.8009,-2.8000) (6.8472,-2.8000) (6.8934,-2.8000) (6.9397,-2.8000) (6.9860,-2.8000) (7.0323,-2.8000) (7.0786,-2.8000) (7.1249,-2.8000) (7.1712,-2.8000) (7.2175,-2.8000) (7.2638,-2.8000) (7.3100,-2.8000) (7.3563,-2.8000) (7.4026,-2.8000) (7.4489,-2.8000) (7.4952,-2.8000) (7.5415,-2.8000) (7.5878,-2.8000) (7.6341,-2.8000) (7.6803,-2.8000) (7.7266,-2.8000) (7.7729,-2.8000) (7.8192,-2.8000) (7.8655,-2.8000) (7.9118,-2.8000) (7.9581,-2.8000) (8.0044,-2.8000) (8.0507,-2.8000) (8.0969,-2.8000) (8.1432,-2.5838) (8.1895,-2.3524) (8.2358,-2.1210) (8.2821,-1.8895) (8.3284,-1.6581) (8.3747,-1.4266) (8.4210,-1.3000) (8.4672,-1.3000) (8.5135,-1.3000) (8.5598,-1.3000) (8.6061,-1.3000) (8.6524,-1.3000) (8.6987,-1.3000) (8.7450,-1.3000) (8.7913,-1.3000) (8.8376,-1.3000) (8.8838,-1.3000) (8.9301,-1.3000) (8.9764,-1.3000) (9.0227,-1.3000) (9.0690,-1.3000) (9.1153,-1.3000) (9.1616,-1.3000) (9.2079,-1.3000) (9.2541,-1.3000) (9.3004,-1.3000) (9.3467,-1.3000) (9.3930,-1.3000) (9.4393,-1.3000) (9.4856,-1.3000) (9.5319,-1.3000) (9.5782,-1.3000) (9.6245,-1.3000) (9.6707,-1.3000) (9.7170,-1.3000) (9.7633,-1.3000) (9.8096,-1.3000) (9.8559,-1.3000) (9.9022,-1.3000) (9.9485,-1.3000) (9.9948,-1.3000) (10.0410,-1.3000) (10.0873,-1.3000) (10.1336,-1.3000) (10.1799,-1.3000) (10.2262,-1.3000) (10.2725,-1.3000) (10.3188,-1.3000) (10.3651,-1.3000) (10.4114,-1.3000) (10.4576,-1.3000) (10.5039,-1.3000) (10.5502,-1.3000) (10.5965,-1.3000) (10.6428,-1.3000) (10.6891,-1.3000) (10.7354,-1.3000) (10.7817,-1.3000) (10.8279,-1.3000) (10.8742,-1.3000) (10.9205,-1.3000) (10.9668,-1.3000) (11.0131,-1.3000) (11.0594,-1.3000) (11.1057,-1.3000) (11.1520,-1.3000) (11.1983,-1.3000) (11.2445,-1.3000) (11.2908,-1.3000) (11.3371,-1.3000) (11.3834,-1.3000) (11.4297,-1.3000) (11.4760,-1.3000) (11.5223,-1.3000) (11.5686,-1.3000) (11.6148,-1.3000) (11.6611,-1.3000) (11.7074,-1.3000) (11.7537,-1.3000) (11.8000,-1.3000)};
\draw[gray,dashed](1.2,-1.3)--(11.8,-1.3);\node[anchor=east] at(1,-1.3){\(1\)};\node[anchor=east] at(1,-2.8){\(0\)};\node at(6,-3.3){归一化后的三个权重};\end{tikzpicture}
\caption[局部覆盖与单位分解拼接]{局部覆盖与单位分解拼接。折线仅示意非负截断的支撑与重叠，实际选取光滑截断。分母在工作区域上为正，归一化后权重和为一；可数情形还需局部有限性。}
\label{b1:fig:visual-f16}
\end{figure}

